% Section 8 — Why it works: representations, transfer and experience (SYNTHETIC results)
\section{What carries the first proposal}
\label{sec:mechanisms}

A first-proposal gain could come from the state encoder, from retrieval in disguise, from the
decoder's prior or from the outcome objective. Each predeclared control removes one of these
(\cref{tab:arms}); \cref{fig:ablation} shows what each costs \synbadge.

\begin{figure}[t]
\centering
% Figure — what carries the gain: component ablations per game, and shared against separate models.
% Reads SYNTHETIC data (gaps EZ-G09, EZ-G11).
\begin{tikzpicture}
\begin{groupplot}[ez,
    group style={group size=2 by 1, horizontal sep=3.1cm},
    height=4.9cm, axis lines*=left, clip=false,
    legend style={font=\sffamily\fontsize{5.6pt}{6.4pt}\selectfont, draw=none, fill=white,
      fill opacity=0.85, text opacity=1, cells={anchor=west}, row sep=-1.5pt},
    error bars/error bar style={line width=0.5pt}, error bars/error mark options={rotate=90, mark size=1.4pt}]
  % ---------------- component ablations
  \nextgroupplot[width=0.44\linewidth, title={Removing one component},
      xmin=-0.125, xmax=0.012, ymin=0.5, ymax=5.5, ytick={1,2,3,4,5}, scaled x ticks=false,
      yticklabels={search-only teacher, retrieval, catalog-bound decoder, imitation only, direct decoder},
      yticklabel style={font=\sffamily\scriptsize\color{ezink}},
      xtick={-0.12,-0.08,-0.04,0}, xticklabel style={/pgf/number format/fixed, /pgf/number format/precision=2},
      xlabel={utility change against full EurekaZero}, legend pos=north west]
    \draw[ezgrey, line width=0.5pt, densely dashed] (axis cs:0,0.5) -- (axis cs:0,5.5);
    \addplot[only marks, mark=*, mark size=1.5pt, ezaccent, ezwhere={domain}{A},
        error bars/.cd, x dir=both, x explicit]
      table[col sep=comma, x=delta, y expr=\thisrow{y}+0.2, x error=err] {data/synthetic/ablations.csv};
    \addlegendentry{A learning procedures}
    \addplot[only marks, mark=square*, mark size=1.3pt, ezgreen, ezwhere={domain}{B},
        error bars/.cd, x dir=both, x explicit]
      table[col sep=comma, x=delta, y=y, x error=err] {data/synthetic/ablations.csv};
    \addlegendentry{B1 next experiment}
    \addplot[only marks, mark=triangle*, mark size=1.7pt, ezwarm, ezwhere={domain}{C},
        error bars/.cd, x dir=both, x explicit]
      table[col sep=comma, x=delta, y expr=\thisrow{y}-0.2, x error=err] {data/synthetic/ablations.csv};
    \addlegendentry{C CAD parts}
  % ---------------- shared against separate models
  \nextgroupplot[width=0.38\linewidth, title={Shared minus separate models},
      xmin=-0.035, xmax=0.025, ymin=0.5, ymax=3.5, ytick={1,2,3}, scaled x ticks=false,
      yticklabels={C CAD parts, B1 next experiment, A learning procedures},
      yticklabel style={font=\sffamily\scriptsize\color{ezink}},
      xtick={-0.03,-0.015,0,0.015}, xticklabel style={/pgf/number format/fixed, /pgf/number format/precision=3},
      xlabel={utility difference}]
    \fill[ezgrey, opacity=0.08] (axis cs:-0.02,0.5) rectangle (axis cs:0.02,3.5);
    \node[font=\sffamily\fontsize{5.6pt}{6.4pt}\selectfont, text=ezgrey, anchor=north] at (axis cs:0,3.48) {equivalence margin $\pm 0.02$};
    \draw[ezgrey, line width=0.5pt, densely dashed] (axis cs:0,0.5) -- (axis cs:0,3.5);
    \addplot[only marks, mark=*, mark size=1.6pt, ezink, error bars/.cd, x dir=both, x explicit]
      table[col sep=comma, x=delta, y=y, x error=err] {data/synthetic/transfer.csv};
\end{groupplot}
\synfigtag
\end{tikzpicture}

\caption{Left: utility lost when one component is removed or replaced, with 95\% intervals, on
development families held out from training. Right: one shared model minus three separate models of
matched total compute; the shaded band is the equivalence margin. Synthetic (gaps EZ-G09, EZ-G11).}
\label{fig:ablation}
\end{figure}

\subsection{The proposer reads the state}

Shuffling states across problems costs \syn[EZ-G09]{\EZshufA}, \syn[EZ-G09]{\EZshufB} and
\syn[EZ-G09]{\EZshufC} utility in the three games, far more than any other ablation and below the
best single solution (\cref{app:outcomes}): the proposer is not emitting one good default. Retrieval of the nearest training
state's incumbent loses \syn[EZ-G09]{\EZabRetA} in learning procedures, so the gain is not a lookup.

\subsection{The solution embedding matters}

Mapping adapter tokens directly to the typed decoder at matched capacity, without the solution
embedding, loses \syn[EZ-G09]{\EZabDirectA} in learning procedures and \syn[EZ-G09]{\EZabDirectC} in
CAD. Two interventions on the embedding itself test whether it carries the move rather than merely
routing information. In \syn[EZ-G09]{\EZswapPairs} pairs of development states with different
incumbents, decoding the embedding of state $s'$ in place of that of $s$ yields a candidate that
scores higher on $s'$ than on $s$ in \syn[EZ-G09]{\EZswapRate\%} [\syn[EZ-G09]{\EZswapRateLo},
\syn[EZ-G09]{\EZswapRateHi}] of pairs, against 50\% if the embedding carried nothing
state-specific; and an embedding taken from a random state of another family beats the recipient's own
first proposal on the recipient's problem in only \syn[EZ-G09]{\EZswapShuffled\%} of pairs. The
decoded candidate follows the embedding, and the embedding follows the state.

\subsection{Open support and outcome training}

A decoder restricted to programs and parts already in the complete-artifact archive loses
\syn[EZ-G09]{\EZabCatA} in learning procedures: generation outside the archive is worth more than
selecting well within it (\cref{prop:closed}). Stopping after conditional imitation (Stage B),
without training on the utility of the first proposal, loses \syn[EZ-G09]{\EZabImitA} and
\syn[EZ-G09]{\EZabImitC}. Imitation teaches the proposer what verified search found; the first-proposal
objective teaches it which of those moves to make when it gets only one.

\subsection{The specification steers the move}

If the proposer reads the objective, changing the objective should change the move. Across
\syn[EZ-G10]{\EZsteerN} reserved families per game we tightened one declared constraint and left
everything else fixed. Tightening the latency limit of a learning procedure produced a faster
committed procedure within the error tolerance in \syn[EZ-G10]{\EZsteerA\%} of families, at a utility
cost of \syn[EZ-G10]{\EZsteerCostA}; tightening the stiffness requirement of a part produced a stiffer
part in \syn[EZ-G10]{\EZsteerC\%} of specifications. A proposer that had memorized good answers per
problem would not respond to an edited specification.

\subsection{Pretrained encoders}

The learning-procedure adapter uses two pretrained tabular foundation encoders whose weights we do not
own (\cref{app:model,app:history}). EurekaZero-Native, which replaces them with encoders trained only
on our corpus, differs by \syn[EZ-G12]{\EZnativeA} [\syn[EZ-G12]{\EZnativeAlo},
\syn[EZ-G12]{\EZnativeAhi}]. The interval excludes a loss larger than 0.022 but crosses the
predeclared margin of 0.02, so non-inferiority is not established: the pretrained encoders may
contribute a small amount, and the headline results do not depend on them for more than that.

\subsection{One model across games}

A shared proposer with three adapters and three decoders, against three separate models of matched
total compute, differs by \syn[EZ-G11]{\EZtrA} [\syn[EZ-G11]{\EZtrAlo}, \syn[EZ-G11]{\EZtrAhi}] in
learning procedures and \syn[EZ-G11]{\EZtrC} [\syn[EZ-G11]{\EZtrClo}, \syn[EZ-G11]{\EZtrChi}] in CAD
(\cref{fig:ablation}). Sharing costs nothing detectable in learning procedures and next-experiment
decisions; in CAD the interval crosses the margin and a small cost cannot be excluded. These results
support one interface for several games. They do not show positive transfer between games, and we do
not claim it (\cref{app:claims}, EZ-C08).

\subsection{Experience and improvement rounds}

\begin{figure}[t]
\centering
% Figure — experience scaling and controlled improvement rounds.
% Reads SYNTHETIC data (gaps EZ-G13, EZ-G14).
\begin{tikzpicture}
\begin{groupplot}[ez,
    group style={group size=2 by 1, horizontal sep=1.6cm},
    width=0.47\linewidth, height=4.7cm, axis lines*=left, clip=false,
    legend style={font=\sffamily\fontsize{5.6pt}{6.4pt}\selectfont, draw=none, fill=white,
      fill opacity=0.85, text opacity=1, cells={anchor=west}, row sep=-1.5pt}]
  % ---------------- experience scaling
  \nextgroupplot[title={More verified experience}, xmode=log, log basis x=2, ymode=log, log basis y=2,
      xmin=0.1, xmax=1.25, ymin=1.5, ymax=100, xtick={0.125,0.25,0.5,1}, xticklabels={1/8,1/4,1/2,1},
      ytick={2,4,8,16,32,64}, yticklabels={2,4,8,16,32,64},
      xlabel={fraction of the training experience}, ylabel={search-equivalent budget $B_{\mathrm{eq}}$},
      legend pos=north west]
    \addplot[ezaccent, line width=1pt, mark=*, mark size=1.5pt] table[col sep=comma, x=fraction, y=beqA] {data/synthetic/history.csv};
    \addlegendentry{A learning procedures (evaluations)}
    \addplot[ezwarm, line width=1pt, mark=triangle*, mark size=1.8pt] table[col sep=comma, x=fraction, y=beqC] {data/synthetic/history.csv};
    \addlegendentry{C CAD parts (simulations)}
  % ---------------- improvement rounds
  \nextgroupplot[title={Controlled improvement rounds (A)}, xmin=-0.3, xmax=3.3, ymin=0, ymax=72,
      xtick={0,1,2,3}, xticklabels={0,1,2,3}, ytick={0,16,32,48,64},
      xlabel={round of training on fresh verified outcomes}, ylabel={$B_{\mathrm{eq}}$ (evaluations)},
      legend pos=north west]
    \addplot[ezaccent, line width=1pt, mark=*, mark size=1.5pt] table[col sep=comma, x=gen, y=ez] {data/synthetic/generations.csv};
    \addlegendentry{teacher seeded by the current model}
    \addplot[ezgrey, line width=0.9pt, densely dashed, mark=o, mark size=1.5pt] table[col sep=comma, x=gen, y=static] {data/synthetic/generations.csv};
    \addlegendentry{static teacher, equal new compute}
\end{groupplot}
\synfigtag
\end{tikzpicture}

\caption{Left: search-equivalent budget against the fraction of verified training experience, on
logarithmic axes. Right: search-equivalent budget after each controlled improvement round, when the
teacher's searches are seeded by the current model, against a static teacher given equal new compute.
Synthetic (gaps EZ-G13, EZ-G14).}
\label{fig:scaling}
\end{figure}

The value of one proposal grows with the verified experience it is trained on. Between one eighth
and all of the training experience, the search-equivalent budget rises from
\syn[EZ-G13]{\EZhistAlow} to \syn[EZ-G13]{\EZbeqA} evaluations in learning procedures and from \syn[EZ-G13]{\EZhistClow} to
\syn[EZ-G13]{\EZbeqC} simulations in CAD (\cref{fig:scaling}), and closed-loop savings rise from
\syn[EZ-G13]{\EZhistSaveLow} to \syn[EZ-G13]{\EZsaveB} times. Training on fresh outcomes from searches seeded by
the current model raises $B_{\mathrm{eq}}$ in each of three rounds, while a static teacher given the
same new compute saturates (\syn[EZ-G14]{\EZstaticThree} against \syn[EZ-G14]{\EZbeqA} evaluations after round 3).

This is the result the program has so far failed to produce. Continued fine-tuning of its first
learned policy on its own accumulated experience raised regression alerts against the frozen first
generation \pre{}; in a matched study of 36 ranking models, training objectives fell throughout while
31 of 36 held-out monitor objectives were worse after 400 updates than after 100 \pre{}; and search
seeded by the frozen policy was worse than retrieval-seeded search on 15 of 17 contexts, at unequal
budgets \pre{} (\cref{app:history}). A rising curve in the right panel, under
the receipt and comparison rules of \cref{app:statistics}, is the single most important gap
in this draft.
